% !TEX root = ../QuantumComputing.tex
%% Lecture 2: Quantum Algorithms
%% Based on original notes by T.J. Osborne, enriched and unified

\section{Quantum algorithms}\label{sec:algorithms}

\begin{historical}[The dawn of quantum algorithms]
  The discovery of efficient quantum algorithms marks one of the great
  intellectual achievements of the late twentieth century.  In 1994, Peter
  Shor~\cite{Shor1994} showed that a quantum computer could factor integers
  in polynomial time, threatening the security of RSA encryption.  Two years
  later, Lov Grover~\cite{Grover1996} demonstrated a quadratic speedup for
  unstructured search.  In 2009, Harrow, Hassidim, and
  Lloyd~\cite{HHL2009} showed that quantum computers could solve certain
  linear systems exponentially faster than any known classical algorithm.
  These three results continue to form the conceptual foundation for most
  quantum algorithms discovered since.
\end{historical}

\subsection{Overview}\label{subsec:algo-overview}

The two most celebrated quantum algorithms are Shor's factoring algorithm
(1994) and Grover's search algorithm (1996).  A third important example,
possibly more relevant for near-term applications, is the
Harrow-Hassidim-Lloyd (HHL) linear-systems algorithm (2009).  For a
comprehensive catalogue of quantum algorithms, see the
\emph{Quantum Algorithm Zoo}\footnote{\url{https://quantumalgorithmzoo.org/}}.

Although there are now very many quantum algorithms, they can often be traced
back fairly directly to one of these three.  It is fair to say that our
understanding of the \emph{discovery} and \emph{design} of quantum
algorithms is still in its infancy.  The field of \emph{quantum software
  engineering}, which would presumably supply robust design practices and
patterns, has not yet been fully developed.

\begin{intuition}[What is quantum advantage?]
  A quantum algorithm provides a \emph{quantum advantage} (or
  \emph{quantum speedup}) when it solves a computational problem using
  asymptotically fewer resources (time, queries, or gates) than any known
  classical algorithm for the same problem.  Advantages can be:
  \begin{itemize}
    \item \textbf{Exponential}: $\bigO{\poly{n}}$ vs.\ $2^{\Omega(n)}$
      (e.g., Shor's algorithm for factoring).
    \item \textbf{Polynomial}: $\bigO{\sqrt{N}}$ vs.\ $\bigO{N}$
      (e.g., Grover's search).
    \item \textbf{Provable}: demonstrated via lower bounds in a query or
      communication model.
    \item \textbf{Heuristic}: observed empirically but without rigorous
      worst-case guarantees (e.g., variational quantum algorithms).
  \end{itemize}
  Understanding which problems admit quantum speedups, and why, remains one
  of the central open questions in quantum computing.
\end{intuition}

\subsection{The query model}\label{subsec:query-model}

To explain the query setting, consider an $N$-bit input
$x = (x_0, \ldots , x_{N-1}) \in \{0, 1\}^{\times N}$.  Usually we take
$N=2^n$, so that bit $x_j$ can be addressed using an $n$-bit index $j$.
One can think of the input as an $N$-bit memory which we can access at any
point of our choice (a ``Random Access Memory'', or RAM).

Memory accesses are modelled via ``black-box'' unitary operations:

\begin{definition}[Standard query oracle]\label{def:standard-oracle}
  The \emph{standard query oracle} for an $N$-bit string
  $x\in\{0,1\}^{\times N}$ is the unitary on $n+1$ qubits defined by
  \begin{equation}
    \eqbox{\Ox:\ket{j,b} \mapsto \ket{j,b\oplus x_j},}
  \end{equation}
  where $j\in \{0,1\}^{\times n}$, $b\in \{0,1\}$, and $\oplus$ denotes
  addition modulo~$2$.  The first $n$ qubits are the \emph{address register}
  and the $(n+1)$st qubit is the \emph{target bit}.
\end{definition}

As a matrix, $\Ox$ is a permutation matrix and hence unitary.  Because $\Ox$
is a quantum operation, a quantum computer can apply it on a
\emph{superposition} of addresses~$j$---something a classical computer
cannot do.

\begin{definition}[Phase oracle]\label{def:phase-oracle}
  Given the ability to make a standard query, we can also make a
  \emph{phase query} by setting the target bit to
  $\keminus = \frac{1}{\sqrt{2}}(\kezero - \keone) = H\keone$:
  \begin{equation}
    \Ox \ket{j}\keminus = (-1)^{x_j}\ket{j}\keminus.
  \end{equation}
  The corresponding $n$-qubit unitary (ignoring the ancilla in the fixed
  state $\keminus$) is
  \begin{equation}
    \eqbox{\Oxpm\ket{j} = (-1)^{x_j}\ket{j}.}
  \end{equation}
\end{definition}

\begin{intuition}[Phase kickback]
  The technique of placing the answer in the \emph{phase} rather than in a
  register is called \emph{phase kickback}.  It is a fundamental trick in
  quantum computing: the ancilla qubit $\keminus$ is an eigenstate of
  $X$ with eigenvalue $-1$, so the oracle's action ``kicks back'' the
  function value $x_j$ as a phase $(-1)^{x_j}$ onto the address register.
  This trick underlies the Deutsch-Jozsa algorithm, Grover's algorithm,
  and quantum phase estimation.
\end{intuition}

\subsection{The Deutsch-Jozsa algorithm}\label{subsec:deutsch-jozsa}

\begin{keyresult}[The Deutsch-Jozsa problem]
  For $N=2^n$, we are given $x\in \{0,1\}^{\times N}$ such that either:
  \begin{enumerate}
    \item all $x_j$ have the same value (``constant''); or
    \item half ($N/2$) of the $x_j$ are $0$ and the other half are $1$
      (``balanced'').
  \end{enumerate}
  The objective is to determine which situation holds.
\end{keyresult}

\begin{algorithmbox}[Deutsch-Jozsa algorithm~\cite{DeutschJozsa1992}]
  \begin{enumerate}
    \item Initialise an $n$-qubit state in $\kezero^{\tp n}$.
    \item Apply a Hadamard gate to each qubit: $H^{\tp n}$.
    \item Apply $\Oxpm$ to the qubits.
    \item Apply $H^{\tp n}$ again.
    \item Measure all qubits in the computational basis.
  \end{enumerate}
  The complete quantum algorithm is the unitary
  \begin{equation}
    \eqbox{U_{\mathrm{DJ}} \equiv H^{\tp n}\Oxpm H^{\tp n}.}
  \end{equation}
\end{algorithmbox}

The circuit for the Deutsch-Jozsa algorithm is:
\begin{center}
\begin{tikzpicture}[scale=0.9]
  % Wires
  \foreach \i in {1,...,3} {
    \draw[qwire] (0,{-\i}) -- (10,{-\i});
  }
  % Dots for n qubits
  \node at (0,-2) {};
  % Labels
  \node[anchor=east,font=\small] at (0,-1) {$\kezero$};
  \node[anchor=east,font=\small] at (0,-2) {$\vdots$};
  \node[anchor=east,font=\small] at (0,-3) {$\kezero$};
  % H gates (first round)
  \node[gate] at (1.5,-1) {$H$};
  \node[gate] at (1.5,-3) {$H$};
  \node at (1.5,-2) {$\vdots$};
  % Oracle
  \node[gate, minimum width=1.2cm, minimum height=2.8cm] at (4,-2) {$\Oxpm$};
  % H gates (second round)
  \node[gate] at (6.5,-1) {$H$};
  \node[gate] at (6.5,-3) {$H$};
  \node at (6.5,-2) {$\vdots$};
  % Measurement
  \qcmeas{8.5}{1}
  \qcmeas{8.5}{3}
  \node at (8.5,-2) {$\vdots$};
\end{tikzpicture}
\end{center}

We follow the state of the quantum computer through each step.  Initially
the state is $\kezero^{\tp n}$.  After the first Hadamard transform:
\begin{equation}
  H^{\tp n}\kezero^{\tp n}
  = \frac{1}{\sqrt{2^n}}\sum_{j\in \{0,1\}^{\times n}} \ket{j}.
\end{equation}
After the phase oracle:
\begin{equation}
  \Oxpm H^{\tp n}\kezero^{\tp n}
  = \frac{1}{\sqrt{2^n}}\sum_{j\in \{0,1\}^{\times n}} (-1)^{x_j}\ket{j}.
\end{equation}
After the second Hadamard:
\begin{equation}
  H^{\tp n}\Oxpm H^{\tp n}\kezero^{\tp n}
  = \frac{1}{2^n}\sum_{j\in \{0,1\}^{\times n}}(-1)^{x_j}
  \sum_{k\in \{0,1\}^{\times n}} (-1)^{j\cdot k}\ket{k},
\end{equation}
where $j\cdot k \equiv \sum_{l=1}^n j_l k_l$.  Since
$j\cdot \mathbf{0} = 0$ for all $j$, the amplitude of the
$\kezero^{\tp n}$ state in the final superposition is
\begin{equation}
  \frac{1}{2^n}\sum_{j\in \{0,1\}^{\times n}}(-1)^{x_j} = \begin{cases}
    1  & \text{if $x_j=0$ for all $j$ (constant $0$),} \\
    -1 & \text{if $x_j=1$ for all $j$ (constant $1$),} \\
    0  & \text{if $x$ is balanced.}
  \end{cases}
\end{equation}
Hence the final measurement yields $\kezero^{\tp n}$ with certainty if $x$
is constant, and yields some other state if $x$ is balanced.

\begin{keyresult}[Deutsch-Jozsa complexity]
  The Deutsch-Jozsa problem can be solved with certainty using only
  \textbf{1~quantum query} and $\bigO{n}$ other operations.  Any classical
  deterministic algorithm needs at least $N/2 + 1$ queries.
\end{keyresult}

\begin{remark}
  The quantum-classical separation for Deutsch-Jozsa only holds for
  \emph{deterministic} algorithms.  A classical \emph{randomised} algorithm
  can solve this problem efficiently: query $x$ at two random positions;
  output ``constant'' if those bits are the same and ``balanced'' if they
  differ.  This succeeds with probability at least $1/2$ when $x$ is
  balanced.
\end{remark}


\subsection{Grover's search algorithm}\label{subsec:grover}

The second most important quantum algorithm after Shor's is Grover's search
algorithm~\cite{Grover1996}.  It provides a \emph{quadratic} speedup---not
exponential---but is far more widely applicable.

\subsubsection{The search problem}

\begin{keyresult}[Unstructured search problem]
  For $N=2^n$, given $x\in \{0,1\}^{\times N}$, find a $j$ such that
  $x_j = 1$ (or output ``no solutions'' if no such $j$ exists).  We denote
  the number of solutions by $t$ (the Hamming weight of $x$).
\end{keyresult}

This problem models searching an $N$-slot unordered database. Classically,
a randomised algorithm needs $\Theta(N)$ queries.  Grover's algorithm solves
it in $\bigO{\sqrt{N}}$ queries and $\bigO{\sqrt{N}\log N}$ other gates.

\subsubsection{The Grover iterate}

We again exploit the phase oracle $\Oxpm\ket{j} = (-1)^{x_j}\ket{j}$.
Define the \emph{zero-reflection}:
\begin{equation}
  R_0\ket{j} \equiv \begin{cases}
    \ket{\mathbf{0}} & \text{if } j= \mathbf{0}, \\
    -\ket{j}         & \text{if } j\not= \mathbf{0}.
  \end{cases}
\end{equation}
That is, $R_0 = 2\ketbra{\mathbf{0}}{\mathbf{0}} - \I$.

\begin{definition}[Grover iterate]\label{def:grover-iterate}
  The \emph{Grover iterate} (or \emph{Grover diffusion operator}) is
  \begin{equation}
    \eqbox{\mathcal{G} \equiv H^{\tp n}R_0 H^{\tp n}\Oxpm
    = \left(2\ket{U}\bra{U} - \I\right)\Oxpm,}
  \end{equation}
  where $\ket{U} = H^{\tp n}\kezero^{\tp n}
  = \frac{1}{\sqrt{N}}\sum_j\ket{j}$ is the uniform superposition.
  A single Grover iterate uses one query and $\bigO{\log N}$ other gates.
\end{definition}

\begin{algorithmbox}[Grover's algorithm~\cite{Grover1996}]
  \begin{enumerate}
    \item Prepare the uniform superposition
      $\ket{U} = H^{\tp n}\kezero^{\tp n}$.
    \item Apply the Grover iterate $\mathcal{G}$ a total of $k$ times (for
      some $k$ to be determined).
    \item Measure the final state in the computational basis.
    \item Verify whether the result $j$ satisfies $x_j=1$.
  \end{enumerate}
\end{algorithmbox}

\subsubsection{Geometric analysis}\label{subsubsec:grover-geometry}

There is an elegant geometric argument for the correctness of Grover's
algorithm.  Define the ``good'' and ``bad'' states:
\begin{equation}
  \ket{G} \equiv  \frac{1}{\sqrt{t}}\sum_{\substack{j:\\x_j=1}}\ket{j},
  \qquad
  \ket{B} \equiv  \frac{1}{\sqrt{N-t}}\sum_{\substack{j:\\x_j=0}}\ket{j}.
\end{equation}
Then the uniform state can be written as
\begin{equation}
  \ket{U} = \sin(\theta)\ket{G} + \cos(\theta)\ket{B},
  \qquad \text{where }
  \theta = \sin^{-1}\!\left(\sqrt{t/N}\right).
\end{equation}

\begin{intuition}[Grover's algorithm as rotation]
  The Grover iterate $\mathcal{G}$ is the product of two reflections in the
  2-dimensional real plane spanned by $\ket{G}$ and $\ket{B}$:
  \begin{enumerate}
    \item $\Oxpm$ reflects through $\ket{B}$
      (i.e., $\Oxpm = \I - 2\ketbra{G}{G}$ restricted to this plane).
    \item $2\ket{U}\bra{U}-\I$ reflects through $\ket{U}$.
  \end{enumerate}
  Each application of $\mathcal{G}$ rotates the state by an angle $2\theta$
  towards $\ket{G}$.

  \begin{center}
    \begin{tikzpicture}[scale=1.8]
      % Axes
      \draw[->,spacecadet,thick] (0,0) -- (2.5,0)
        node[right,font=\small] {$\ket{B}$};
      \draw[->,spacecadet,thick] (0,0) -- (0,2.2)
        node[above,font=\small] {$\ket{G}$};
      % Initial state |U>
      \draw[->,cgblue,very thick] (0,0) -- ({2*cos(20)},{2*sin(20)})
        node[above right,font=\small] {$\ket{U}$};
      % After 1 iterate
      \draw[->,munsell,very thick] (0,0) -- ({2*cos(60)},{2*sin(60)})
        node[above right,font=\small] {$\mathcal{G}\ket{U}$};
      % After 2 iterates
      \draw[->,banana!80!black,very thick,dashed] (0,0) -- ({2*cos(100)},{2*sin(100)})
        node[above left,font=\small] {$\mathcal{G}^2\ket{U}$};
      % Angle arcs
      \draw[cgblue,thin] (0.8,0) arc (0:20:0.8);
      \node[cgblue,font=\scriptsize] at (1.0,0.15) {$\theta$};
      \draw[munsell,thin] (0.6,0) arc (0:60:0.6);
      \node[munsell,font=\scriptsize] at (0.55,0.35) {$3\theta$};
      % Unit circle arc
      \draw[spacecadet!30,thin,dashed] (2,0) arc (0:110:2);
    \end{tikzpicture}
  \end{center}
\end{intuition}

After $k$ applications of $\mathcal{G}$, the state becomes
\begin{equation}
  \mathcal{G}^k\ket{U}
  = \sin\bigl((2k + 1)\theta\bigr)\ket{G}
  + \cos\bigl((2k + 1)\theta\bigr)\ket{B}.
\end{equation}
The probability of observing a solution upon measurement is
\begin{equation}
  P_k = \sin^2\bigl((2k + 1)\theta\bigr).
\end{equation}

We want $P_k$ close to $1$.  Setting
$\widetilde{k} = \frac{\pi}{4\theta} - \frac{1}{2}$ gives
$(2\widetilde{k}+1)\theta = \pi/2$ and hence
$P_{\widetilde{k}} = 1$.

\begin{example}[$t = N/4$]
  If $t = N/4$ then $\theta = \pi/6$ and $\widetilde{k} = 1$.
  Thus a single Grover iteration suffices to find a solution with certainty.
\end{example}

In general $\widetilde{k}$ is not an integer.  Choosing $k$ to be the
nearest integer, we bound the failure probability:
\begin{equation}
  \begin{split}
    1-P_k &= \cos^2\bigl((2k+1)\theta\bigr)
    = \sin^2\bigl(2(k-\widetilde{k})\theta\bigr) \\
    &\le \sin^2(\theta) = \frac{t}{N},
  \end{split}
\end{equation}
where we used $|k-\widetilde{k}|\le \frac{1}{2}$.  Since
$\sin^{-1}(\sqrt{t/N})\ge \sqrt{t/N}$, the number of queries is
\begin{equation}
  \eqbox{k\le \frac{\pi}{4}\sqrt{\frac{N}{t}}.}
\end{equation}

\subsubsection{Variants and extensions of Grover's algorithm}

\begin{itemize}
  \item \textbf{Exact search}: If $t$ is known exactly, the algorithm can
    be tweaked to end up in exactly $\ket{G}$ by adjusting $\theta$.

  \item \textbf{Unknown $t$}: If we do not know $t$, running the algorithm
    with exponentially increasing guesses for $k$ still gives an expected
    query count of $\bigO{\sqrt{N/t}}$.  If $t=0$, this is detected by
    verifying $x_j$ for the output $j$.

  \item \textbf{Known lower bound}: If $\tau \le t$ is a known lower bound,
    the algorithm uses an expected $\bigO{\sqrt{N/\tau}}$ queries.

  \item \textbf{Boosted error reduction}: To reduce the failure probability
    to $\epsilon$, one can use $\bigO{\sqrt{N \log(1/\epsilon)}}$
    queries---note the $\log(1/\epsilon)$ is \emph{inside} the square root,
    which is better than naive repetition giving
    $\bigO{\sqrt{N}\log(1/\epsilon)}$.
\end{itemize}

\begin{remark}
  Grover's $\bigO{\sqrt{N}}$ query complexity is \emph{optimal}: Bennett,
  Bernstein, Brassard, and Vazirani (1997) proved that any quantum algorithm
  for unstructured search requires $\Omega(\sqrt{N})$ queries.  This is one
  of the few known tight quantum lower bounds.
\end{remark}
